\documentclass[a4paper,11pt]{article}
\usepackage[T2A]{fontenc}
\usepackage[utf8]{inputenc}
\usepackage[russian]{babel}
\usepackage{amsmath,amssymb,mathtools}
\usepackage{geometry}
\geometry{left=23mm,right=23mm,top=24mm,bottom=24mm}
\usepackage{microtype}
\usepackage{icomma}
\usepackage{tikz}
\usetikzlibrary{calc,arrows.meta,positioning,shapes.geometric}
\usepackage{tkz-euclide}
\usepackage{booktabs,tabularx,array}
\usepackage{enumitem}
\usepackage{fancyhdr}
\usepackage{hyperref}
\usepackage{xcolor}
\usepackage{setspace}
\definecolor{accent}{HTML}{008080}
\definecolor{darkaccent}{HTML}{075554}
\definecolor{ink}{HTML}{283338}
\definecolor{muted}{HTML}{607073}
\hypersetup{colorlinks=true,linkcolor=darkaccent,urlcolor=accent,pdftitle={Чек-лист. Прямые призмы и правильные треугольные пирамиды},pdfauthor={Лёвин Артём Александрович}}
\setlength{\parindent}{0pt}
\setlength{\parskip}{5pt}
\setstretch{1.16}
\setlist[enumerate]{leftmargin=*,itemsep=5pt,topsep=3pt}
\setlist[itemize]{leftmargin=*,itemsep=2pt,topsep=2pt}
\pagestyle{fancy}
\fancyhf{}
\fancyhead[L]{\small\color{muted}Призмы и пирамиды}
\fancyhead[R]{\small\color{muted}Николь Саркисьянц}
\fancyfoot[C]{\small\color{muted}\thepage}
\renewcommand{\headrulewidth}{0.25pt}
\renewcommand{\footrulewidth}{0pt}
\setlength{\headheight}{15pt}
\newcommand{\topic}[1]{\vspace{6pt}{\Large\bfseries\color{darkaccent}#1}\par\vspace{4pt}\hrule\vspace{12pt}}
\newcommand{\point}[1]{\textcolor{darkaccent}{\textbf{#1}}}
\newcommand{\smhead}[1]{\medskip\textbf{\color{darkaccent}#1}\par}
\newcommand{\ans}[1]{\textcolor{darkaccent}{\textbf{Ответ:}}\ $#1$.}
\begin{document}
\thispagestyle{empty}
\vspace*{28mm}
\begin{center}
{\large\color{darkaccent}МАТЕМАТИКА}\par
\vspace{26mm}
{\Huge Чек-лист}\par
\vspace{13mm}
{\LARGE\color{accent}Прямые призмы\\[3pt]и правильные треугольные пирамиды}\par
\vspace{18mm}
{\large Подготовка к ЕГЭ по профильной математике}\par
\vspace{37mm}
{\Large Николь Саркисьянц}\par
\vspace{11mm}
{\large 11 ноября 2026 г.}\par
\vfill
{\normalsize Лёвин Артём Александрович}\\[3pt]
{\normalsize эксклюзивно для Николь Саркисьянц}\par
\vspace{15mm}
\begin{tikzpicture}[x=1cm,y=1cm]
\draw[accent,line width=0.8pt] (-4.0,0) .. controls (-2.6,0.75) and (-1.7,-0.70) .. (0,0)
                                  .. controls (1.4,0.67) and (2.8,-0.65) .. (4,0);
\fill[darkaccent] (0,0) circle (1.5pt);
\end{tikzpicture}
\end{center}
\clearpage

\topic{1. Прямая призма: площадь поверхности}
\point{Главное свойство.} У прямой призмы боковые рёбра перпендикулярны основанию, поэтому длина бокового ребра равна высоте $H$.
\[
\boxed{S_{\text{бок}}=P_{\text{осн}}H},\qquad
\boxed{S_{\text{полн}}=S_{\text{бок}}+2S_{\text{осн}}}.
\]
Здесь $P_{\text{осн}}$ --- периметр основания, $S_{\text{осн}}$ --- площадь \emph{одного} основания. Если в условии просто сказано «площадь поверхности», ищем полную площадь.

\smhead{Ромб в основании}
Диагонали ромба взаимно перпендикулярны и делятся точкой пересечения пополам. Если диагонали $d_1$ и $d_2$, то
\[
S_{\text{осн}}=\frac{d_1d_2}{2},\qquad
 a=\sqrt{\left(\frac{d_1}{2}\right)^2+\left(\frac{d_2}{2}\right)^2},
 \qquad P_{\text{осн}}=4a.
\]
\begin{center}
\begin{tikzpicture}[scale=0.66]
\tkzDefPoint(-1.5,0){A}\tkzDefPoint(0,2){B}\tkzDefPoint(1.5,0){C}\tkzDefPoint(0,-2){D}\tkzDefPoint(0,0){O}
\tkzDrawPolygon[color=darkaccent,thick](A,B,C,D)
\tkzDrawSegments[dashed,color=accent](A,C B,D)
\tkzDrawPoints[size=2](A,B,C,D,O)
\tkzLabelPoints[left](A)\tkzLabelPoints[right](C)
\tkzLabelPoints[above](B)\tkzLabelPoints[below](D)
\tkzLabelPoint[above right](O){$O$}
\end{tikzpicture}
\quad\begin{minipage}{0.47\linewidth}
\small
\textbf{Ключевой прямоугольный треугольник:} $AOD$.\\[4pt]
$AO=\frac{AC}{2}$, $DO=\frac{BD}{2}$,\\[4pt]
$AD^2=AO^2+DO^2$.
\end{minipage}
\end{center}
\smhead{Пример 1. Площадь полной поверхности}
В основании прямой призмы ромб с диагоналями $3$ и $4$; боковое ребро равно $5$.
\[
S_{\text{осн}}=\frac{3\cdot4}{2}=6,\qquad
 a=\sqrt{1{,}5^2+2^2}=2{,}5.
\]
\[
S_{\text{бок}}=4\cdot2{,}5\cdot5=50,\qquad
S_{\text{полн}}=50+2\cdot6=\boxed{62}.
\]
\point{Проверка:} в полной площади обязательно учтены два основания; в боковой --- только боковые грани.

\clearpage
\topic{2. Обратная задача и сравнение призм}
\smhead{Пример 2. По полной площади найти высоту}
Основание --- ромб со стороной $5$ и площадью $24$. Полная площадь поверхности прямой призмы равна $248$. Пусть её высота равна $x$.
\[
P_{\text{осн}}=4\cdot5=20,\qquad
248=20x+2\cdot24.
\]
\[
20x=248-48=200,\qquad \boxed{x=10}.
\]
\point{Внимание:} $248$ --- полная площадь, а $20x$ --- боковая. Их нельзя приравнивать без учёта оснований.

\smhead{Метод 1. Отношение одинаковых величин}
Если две прямые призмы имеют одну и ту же высоту $H$, то
\[
\frac{S_{\text{бок, мал}}}{S_{\text{бок, бол}}}
=\frac{P_{\text{мал}}H}{P_{\text{бол}}H}
=\frac{P_{\text{мал}}}{P_{\text{бол}}}.
\]
Если основание малой треугольной призмы образовано средними линиями исходного треугольника, её стороны вдвое меньше соответствующих сторон исходного основания. Значит, периметр и боковая площадь также уменьшаются вдвое.

\textbf{Пример 3.} Если боковая площадь исходной призмы равна $12$, то
\[
S_{\text{бок, мал}}=12\cdot\frac12=\boxed{6}.
\]

\smhead{Метод 2. Множители изменения}
Для прямой призмы $V=S_{\text{осн}}H$. У равностороннего треугольника $S_{\text{осн}}=\frac{a^2\sqrt3}{4}$. Поэтому при изменении стороны $a$ в $k$ раз и высоты $H$ в $m$ раз объём изменяется в $k^2m$ раз:
\[
\boxed{\frac{V_2}{V_1}=\left(\frac{a_2}{a_1}\right)^2\frac{H_2}{H_1}}.
\]
\textbf{Пример 4.} Объём первоначальной правильной треугольной призмы равен $6$. Сторону основания увеличили втрое, высоту уменьшили вдвое:
\[
V_2=6\cdot3^2\cdot\frac12=\boxed{27}.
\]
\point{Оба метода нужно освоить:} сначала составить отношение формул и сократить общие множители; затем проверить результат быстрым анализом изменения каждой величины.

\clearpage
\topic{Алгоритм: сравнение геометрических величин}
\vspace{9mm}
\begin{center}
\begin{tikzpicture}[
  every node/.style={font=\small,align=center},
  start/.style={ellipse,draw=darkaccent,thick,minimum width=6.1cm,minimum height=1.05cm},
  flowstep/.style={rectangle,rounded corners=2pt,draw=accent,thick,text width=9.5cm,minimum height=1.25cm,inner sep=8pt},
  arrow/.style={-{Stealth[length=3mm]},thick,color=darkaccent},
  node distance=12mm]
\node[start] (n0) {Даны две фигуры и величина у одной из них};
\node[flowstep,below=of n0] (n1) {Запишите формулу искомой величины\\для обеих фигур: $S=P H$ или $V=S_{\text{осн}}H$};
\node[flowstep,below=of n1] (n2) {Разделите вторую формулу на первую,\\сократите общие множители};
\node[flowstep,below=of n2] (n3) {Выразите оставшиеся отношения\\через изменения сторон и высот};
\node[flowstep,below=of n3] (n4) {Подставьте известную величину,\\найдите искомую и проверьте результат};
\node[start,below=of n4] (n5) {Сравните с экспресс-методом множителей};
\draw[arrow] (n0)--(n1);
\draw[arrow] (n1)--(n2);
\draw[arrow] (n2)--(n3);
\draw[arrow] (n3)--(n4);
\draw[arrow] (n4)--(n5);
\end{tikzpicture}
\end{center}
\clearpage

\topic{3. Правильная треугольная пирамида}
Пусть $SABC$ --- правильная треугольная пирамида. Тогда $ABC$ --- равносторонний треугольник; боковые рёбра $SA=SB=SC=b$, стороны основания $AB=BC=CA=a$; вершина $S$ расположена над центром $O$ основания.
\[
\boxed{S_{\text{осн}}=\frac{a^2\sqrt3}{4}},\qquad
\boxed{V=\frac13 S_{\text{осн}}H},\qquad
\boxed{S_{\text{полн}}=S_{\text{бок}}+S_{\text{осн}}}.
\]
\textbf{Обозначения:} $H=SO$ --- высота пирамиды, перпендикулярная плоскости основания; $h=SM$ --- апофема (высота боковой грани), $M$ --- середина $BC$.
\begin{center}
\begin{tikzpicture}[scale=0.85]
\tkzDefPoint(-2.5,-0.7){A}\tkzDefPoint(2.7,-0.7){B}\tkzDefPoint(0.3,1.2){C}
\tkzDefPoint(1.5,0.25){M}\tkzDefPoint(0.1667,-0.0667){O}\tkzDefPoint(0.1667,3.15){S}
\tkzDrawSegments[color=darkaccent,thick](A,B B,C C,A S,A S,B S,C)
\tkzDrawSegments[dashed,color=muted](A,M S,O)
\tkzDrawSegments[color=accent,thick](S,M)
\tkzDrawPoints[size=2](A,B,C,S,M,O)
\tkzLabelPoints[below left](A)\tkzLabelPoints[below right](B)
\tkzLabelPoints[right](C)\tkzLabelPoints[above](S)
\tkzLabelPoint[right](M){$M$}\tkzLabelPoint[left](O){$O$}
\end{tikzpicture}
\quad
\begin{minipage}{0.43\linewidth}
\small
\textbf{Разные высоты:}\\[5pt]
$SO\perp(ABC)$ --- высота всей пирамиды.\\[5pt]
$SM\perp BC$ --- высота боковой грани.\\[5pt]
$AM$ --- медиана основания; $O\in AM$.
\end{minipage}
\end{center}
Каждая боковая грань --- равнобедренный треугольник с основанием $a$ и высотой $h$, поэтому
\[
\boxed{S_{\text{бок}}=3\cdot\frac{ah}{2}=\frac12 P_{\text{осн}}h=ph},
\qquad p=\frac{P_{\text{осн}}}{2}=\frac{3a}{2}.
\]
Формула $S_{\text{бок}}=ph$ пригодна и для других пирамид, \emph{если высоты всех боковых граней равны} и $h$ обозначает их общую высоту.

\clearpage
\topic{4. Как найти высоту и апофему}
\smhead{Построение 1. Центр основания}
В равностороннем $\triangle ABC$ медианы одновременно служат высотами и биссектрисами. Они пересекаются в точке $O$ и делятся ею в отношении $2:1$ от вершины. Для медианы $AM$:
\[
AM=\sqrt{a^2-\left(\frac a2\right)^2}=\frac{a\sqrt3}{2},
\qquad AO=\frac23 AM=\frac{a\sqrt3}{3}.
\]
\smhead{Построение 2. Высота всей пирамиды}
В прямоугольном треугольнике $SAO$ гипотенуза $SA=b$, а $SO=H$:
\[
SA^2=SO^2+AO^2\quad\Longrightarrow\quad
\boxed{H=SO=\sqrt{b^2-\frac{a^2}{3}}}.
\]
\smhead{Построение 3. Апофема боковой грани}
Треугольник $SBM$ прямоугольный, поскольку $M$ --- середина основания равнобедренной грани $SBC$. Поэтому
\[
SM^2=SB^2-BM^2=b^2-\left(\frac a2\right)^2,
\qquad \boxed{h=SM=\sqrt{b^2-\frac{a^2}{4}}}.
\]
\point{Условие существования невырожденной пирамиды:} $a>0$, $b>0$, $b^2>\frac{a^2}{3}$. Перед применением формул определите, какую высоту требуется найти: $SO$ или $SM$.

\smhead{Схема выбора прямоугольного треугольника}
\begin{tabularx}{\linewidth}{@{}>{\bfseries}p{0.30\linewidth}X@{}}
\toprule
Нужно найти & Какой треугольник использовать \\
\midrule
Медиану $AM$ & $\triangle ABM$ (или $\triangle ACM$) в основании \\
Отрезок $AO$ & $AO=\dfrac23 AM$ \\
Высоту $SO$ & $\triangle SAO$ \\
Апофему $SM$ & $\triangle SBM$ \\
\bottomrule
\end{tabularx}

\smhead{Ошибки, которые стоит исключить}
\begin{itemize}
\item У призмы в полной площади \textbf{два} основания, у пирамиды --- \textbf{одно}.
\item У пирамиды объём содержит множитель $\frac13$; у призмы его нет.
\item В формуле боковой площади пирамиды используется апофема $h$, в формуле объёма --- высота $H$.
\item При увеличении стороны равностороннего основания в $k$ раз его площадь меняется в $k^2$ раз.
\item Правильная пирамида не обязана быть правильным тетраэдром: боковое ребро $b$ вообще говоря отличается от $a$.
\end{itemize}

\clearpage
\topic{5. Разобранные задачи на пирамиду}
\smhead{Пример 5. Известна сторона основания и высота}
В правильной треугольной пирамиде $a=1$, высота пирамиды $H=\sqrt3$. Найти объём.
\[
S_{\text{осн}}=\frac{1^2\sqrt3}{4}=\frac{\sqrt3}{4},\qquad
V=\frac13\cdot\frac{\sqrt3}{4}\cdot\sqrt3=\boxed{\frac14}.
\]
\point{Ключ к чтению условия:} если сказано просто «высота пирамиды», имеется в виду перпендикуляр к плоскости основания, то есть $SO$.

\smhead{Пример 6. Боковые рёбра попарно перпендикулярны}
Дано: $SA=SB=SC=3$, углы между каждой парой боковых рёбер при вершине $S$ прямые. Найти объём пирамиды.

Из прямоугольного треугольника $ASB$:
\[
AB=\sqrt{SA^2+SB^2}=3\sqrt2.
\]
Аналогично $BC=AC=3\sqrt2$: основание равностороннее. Значит, формулы для правильной треугольной пирамиды применимы.
\[
S_{\text{осн}}=\frac{(3\sqrt2)^2\sqrt3}{4}=\frac{9\sqrt3}{2},
\qquad H=\sqrt{3^2-\frac{(3\sqrt2)^2}{3}}=\sqrt3.
\]
\[
V=\frac13\cdot\frac{9\sqrt3}{2}\cdot\sqrt3=\boxed{\frac92}.
\]
\smhead{Пример 7. Высота как отрезок до пересечения медиан}
Пирамида правильная, площадь основания равна $2$, объём $6$. Точка $O$ --- пересечение биссектрис основания. Найти $SO$.

В равностороннем треугольнике точка $O$ --- также пересечение медиан и основание высоты пирамиды. Поэтому
\[
6=\frac13\cdot2\cdot SO\quad\Longrightarrow\quad \boxed{SO=9}.
\]
\point{Экспресс-проверка:} подставьте найденную высоту обратно в формулу объёма. Длину стороны основания искать здесь необязательно.

\clearpage
\topic{Тренировочный блок}
\textcolor{muted}{Путь А. 10 задач: от базовых формул к комбинированным построениям. В задачах 4 и 5 используйте два способа: отношение формул и множители изменения.}
\begin{enumerate}[label=\textbf{\arabic*.},itemsep=11pt]
\item В основании прямой призмы ромб с диагоналями $6$ и $8$. Высота призмы равна $7$. Найдите площадь полной поверхности.
\item В основании прямой призмы ромб со стороной $13$ и диагоналями $10$ и $24$. Площадь полной поверхности равна $500$. Найдите высоту призмы.
\item Боковая площадь прямой треугольной призмы равна $30$. Через средние линии основания отсечена прямая треугольная призма той же высоты. Найдите боковую площадь отсечённой призмы.
\item Объём правильной треугольной призмы равен $8$. Сторону основания увеличили вдвое, высоту --- втрое. Найдите новый объём.
\item Объём правильной треугольной призмы равен $18$. Сторону основания увеличили втрое, а высоту уменьшили втрое. Найдите новый объём.
\item Сторона основания правильной треугольной пирамиды равна $4$, а её высота --- $3\sqrt3$. Найдите объём.
\item В правильной треугольной пирамиде сторона основания равна $6$, боковое ребро --- $5$. Найдите площадь полной поверхности.
\item Три боковых ребра треугольной пирамиды, выходящие из одной вершины, попарно перпендикулярны; каждое ребро равно $3$. Найдите объём пирамиды.
\item Площадь основания правильной треугольной пирамиды равна $8$, а объём --- $16$. Найдите высоту пирамиды.
\item В правильной треугольной пирамиде сторона основания равна $6$, высота пирамиды равна $4$. Найдите высоту боковой грани и площадь боковой поверхности.
\end{enumerate}
\clearpage
\topic{Ответы}
\begin{center}
\renewcommand{\arraystretch}{1.35}
\begin{tabularx}{0.88\linewidth}{@{}>{\bfseries}p{0.12\linewidth}X@{}}
\toprule
№ & Ответ \\
\midrule
1 & $188$ \\
2 & $5$ \\
3 & $15$ \\
4 & $96$ \\
5 & $54$ \\
6 & $12$ \\
7 & $36+9\sqrt3$ \\
8 & $\dfrac92$ \\
9 & $6$ \\
10 & $\sqrt{19}$; $9\sqrt{19}$ \\
\bottomrule
\end{tabularx}
\end{center}
\vspace{13mm}
\smhead{Итоговая памятка}
\begin{tabularx}{\linewidth}{@{}>{\bfseries}p{0.39\linewidth}X@{}}
\toprule
Сигнал в условии & Что делать \\
\midrule
Прямая призма, поверхность & $S_{\text{полн}}=P_{\text{осн}}H+2S_{\text{осн}}$ \\
Ромб и диагонали & $S=\dfrac{d_1d_2}{2}$; сторону найти по половинам диагоналей \\
Одинаковая высота призм & Сравнить периметры для боковых площадей \\
Изменили сторону и высоту & Учесть квадрат стороны в площади основания \\
Правильная треугольная пирамида & $V=\dfrac13S_{\text{осн}}H$ \\
Нужна высота $SO$ & Медиана основания $\to AO\to SO$ \\
Нужна апофема $SM$ & Прямоугольный треугольник боковой грани \\
\bottomrule
\end{tabularx}
\vfill
\begin{center}\small\color{muted}
Лёвин Артём Александрович\\[-1pt]
эксклюзивно для Николь Саркисьянц
\end{center}
\end{document}
